% Created 2026-09-15 Tue 18:44
% Intended LaTeX compiler: pdflatex
\documentclass[11pt]{article}
\usepackage[utf8]{inputenc}
\usepackage[T1]{fontenc}
\usepackage{graphicx}
\usepackage{longtable}
\usepackage{wrapfig}
\usepackage{rotating}
\usepackage[normalem]{ulem}
\usepackage{amsmath}
\usepackage{amssymb}
\usepackage{capt-of}
\usepackage{hyperref}
\author{fcalado}
\date{\today}
\title{}
\hypersetup{
 pdfauthor={fcalado},
 pdftitle={},
 pdfkeywords={},
 pdfsubject={},
 pdfcreator={Emacs 29.3 (Org mode 9.6.15)}, 
 pdflang={English}}
\begin{document}

\tableofcontents

\section{draft}
\label{sec:org39c6bfd}
\subsection{introduction}
\label{sec:org1a815bb}
This teaching statement covers my teaching philosophy, the courses and
workshops I've taught at Pratt, mentoring and advising, and future
plans. 

\subsection{teaching philosophy}
\label{sec:orga9c8223}
My teaching philosophy draws from my graduate school training as a
literary studies and queer studies scholar. Many of the methods here
were developed in my role as a writing instructor for 6 years across
three different institutions.

This philosophy is organized into four major themes: (1) close reading
as critical and creative practice; (2) writing as thinking; (3) all
knowledge as contingent; and (4) active \& experiential learning.

\subsubsection{1) close reading as critical practice:}
\label{sec:org31671de}
The most important skill that I teach in my classes is close reading.
Originally developed as an analytical method for literary criticism,
close reading is defined as the close and sustained attention to a
selection of evidence, such as a passage in a text. It emphasizes the
details of the evidence, such as individual words and sentence
structure, over the general. While close reading is the analytical
methodology par excelance for Literary Studies, as well as other
fields the Humanities, its principles apply across different
analytical contexts, particularly across my Pratt SI courses, from
INFO 601: Foundations to INFO 664: Programming Cultural Heritage.

For example, in my 601: Foundations class, we often use close reading
as a tool to guide discussion on a specific reading. I direct students
to identify passages that struck them as provocative or confusing,
which we then read together as a class. The opportunity to put a
single passage on the projector proves generative, giving students the
opportunity to dig into details and make new connections in real time.
In my \url{664} class, I take a similar approach. For example, I
encourage my students to consider even large datasets at a close
distance. I use my own research as an example to model how close
attention to language offers material for critique. Last year, I had
several students develop 664 projects which cusomized Small Language
Models (SLMs) in order to better study language data up close, with
one group of students from my course ended up winning the "People's
Choice Award" at the 2026 Info Show.

\subsubsection{2) writing as active learning:}
\label{sec:orgdd5ed44}
Another technique that I often use in courses is writing. This is an
active learning technique which I employ through formal writing
assignments, which are highly scaffolded papers developed over several
weeks, and low-stakes writing exercises, which is a daily practice for
organizing thoughts and generating ideas.

In my discussion-based courses, I use free-writing assignments often
as a precursor to class discussion. However, I also use free-writing
in my programming class, INFO 607: Special Topics (How to Build a
Bot), as a tool to help students theorize their final projects. Formal
writing assignments, by contrast, are highly scaffolded, with detailed
guidance on each step, from drafting research questions, to outlining
arguments and evidence, to crafting paragraphs. In my courses that
require formal writing assignments, such \url{601}, INFO: 657, and
INFO: 656, I dedicate time across several class meetings to addressing
individual aspects of the paper, such as the "hook" of a typical
research paper, or the "Discussion" section of a conference paper.
This hands-on practice with writing in class is even more important
now, with the rise of automated tools which tempt students to offload
writing tasks, and therefore the development of critical thinking
skills.

\subsubsection{3) knowledge as contingent:}
\label{sec:org1ad0ba4}
I maintain the importance of meeting students where they are by
inviting them to use their own life stories and personal interests as
a ground for critique and understanding. I also emphasize the
importance of perspective, that all knowledge is contingent and that
"infinite vision" as Donna Haraway asserts, "is an illusion, a god
trick." Students thus learn the contingency of all knowledge through
exposure to a wide range of theoretical points of view and by drawing
from their own life experience.

I employ these principles through the selection of readings in my
courses, particularly INFO 601, INFO 656, and INFO 657. Readings range
from different disciplines and fields, showing how differences in
initial questions, assumptions, and approaches will lead to quite
different conclusions. For example, in INFO: 657 I assign readings
from different fields in the humanities (such as Black Feminist
Studies and Archival Studies) so that students have a basis for
understanding how the field of Digital Humanities departs from others.

\subsection{courses \& workshops taught}
\label{sec:org60ec2d0}
Courses
\begin{itemize}
\item INFO 656: Machine Learning, Fall 2026.
\begin{itemize}
\item This course offers an introduction to machine learning as a
practical tool that we can use, and as a technological field with
social implications. We will learn about key concepts in machine
learning; survey a few key machine learning techniques, such as
supervised methods for machine learning (regression and
classification), which attempt to map data onto desired outputs,
and unsupervised methods (clustering and association), which
attempt to find structure within data itself; use openly available
tools to implement these techniques on text and image data; learn
how to assess the effectiveness of different techniques on
particular datasets; and discuss basic issues that confront all
machine learning methods.
\end{itemize}
\item INFO 657: Digital Humanities. Spring 2026.
\begin{itemize}
\item This course examines the history, theory, and practice of digital
humanities, paying special attention to the ways in which digital
humanities are transforming research, disciplines, and even the
academy itself. Topics include contrasts and continuities between
traditional and digital humanities; tools and techniques used by
digital humanists; the processes of planning, funding, managing,
and evaluating digital humanities projects; ways in which the
digital humanities impact scholarly communication and higher
education; and the special roles of libraries and information
professionals in this growing movement.
\end{itemize}
\item INFO 697: Special Topics - How to Build a Bot. Summer 2025.
\begin{itemize}
\item This course offers a practical introduction to building bots in
Python alongside a critical examination of algorithmic bias.
Students will learn core programming skills like data collection
and API interaction in order to create web crawlers and social
media bots. In parallel, students will explore the ethical
consequences of automated systems on social platforms, such as the
amplification of misinformation and bias that perpetuate social
inequalities and discrimination. This course will equip students
with the critical perspectives and technical skills to analyze
automated systems in a world increasingly shaped by AI and
algorithmic decision-making. This course is intended for students
interested in both technical development and the social impacts of
automation.
\end{itemize}
\item INFO 664: Programming Cultural Heritage. Fall 2024, Spring 2025,
Fall 2025, Spring 2026, Fall 2026.
\begin{itemize}
\item Technical knowledge and proficiency is an increasingly valued and
useful skill in a cultural heritage information professional’s
skillset. This course will introduce and provide a platform for
the building up of these abilities. Over the course of the
semester the student will become familiar with programming,
working with data in a programmatic manner and applying these
skills to build a demonstrable project. This course will focus on
teaching the basics that will enable an information professional
to be successful in today’s digital world.
\end{itemize}
\item INFO 601: Foundations of Information. Fall 2024, Spring 2025,
Fall 2025, Spring 2026, Fall 2026.
\begin{itemize}
\item This foundational course focuses on the intersection of people,
information, and technology and the theoretical and conceptual
foundations of the information field. Students will be introduced
to ideas and concepts that will inform future specializations in
their course of study and provide them with concrete strategies
for ongoing professional growth and development in their area of
interest.
\end{itemize}
\end{itemize}

Workshops
\begin{itemize}
\item Text Processing with Python: Exploring \emph{The Odyssey}. \emph{The Data
Storytelling Lab} (\emph{DSL}). Fall 2025, Fall 2026.
\begin{itemize}
\item In this hands-on Python workshop, you’ll learn the fundamentals of
text processing using spaCy (one of the most widely used Python
tools for working with text) to study different versions of
Homer’s epic poem, The Odyssey. We will use spaCy to tokenize the
different versions of the text, identify persons and places, and
compare linguistic patterns across adaptations---allowing us to
see the subtle differences in style, characterization, and
narrative emphasis across the story’s different versions. How does
The Odyssey change across the major translations? What are the
differences between Emily Wilson’s version (the one used as the
basis for Christopher Nolan’s recent feature film) and other major
versions of the Odyssey? Come and explore the narrative features
of the text for yourself.s
\end{itemize}
\end{itemize}


\subsection{mentoring \& advising}
\label{sec:org4cae90e}
In addition to the \textasciitilde{}30 students that I teach directly each semester, I
currently advise 35 students, mostly in Library and Information
Science (LIS), but also a few in Data Analytics and Visualization
(DAV), and Information Experience Design (IXD). Each semester, I meet
individually with many of these students to discuss their progress in
coursework and projects and to strategize their final portfolios,
which is the capstone project for the program. I also discuss career
expectations and goals with students.

In addition to advising students individually, I also advise two
students groups. For one group, Latines in Tech, I am the official
Faculty Advisor. Latines in Tech is professional development group
primarily focused on Latinx students. Their main goals are to provide
professional development and networking opportunities for their
members. The second group, The \emph{Stacks} Student Journal, which is
still in the process of obtaining official student group status, is an
editorial board comprised of some former students from my 601:
Foundations course. These students are currently producing a
peer-reviewed journal to showcase their 601 papers. My support for
these groups invoves attending meetings, participating in their
events, and offering guidance as needed.


\subsection{future plans}
\label{sec:org81bf273}
My future plans involve developing two courses:
\begin{itemize}
\item First, I plan to develop a Python programming course that is
specific to working with text-based data. This course would build
off my workshop that I run for the Data Storytelling Lab, which
focuses on using advanced tools like the \emph{spaCy} library in Python,
which allows users to perform advanced tasks in text cleaning,
processing, annotating, and NLP analysis, like Named Entity
Recognition. This course would join my other course, INFO 697: How
to Build a Bot, to offer instruction in intermediate programming
skills.
\item Second, I plan to develop a "Readings in Queer Studies" course.
While I am aware that theoretically heavy courses have not
historically done well in this program, I think attitudes and
appetites for this material is shifting. When assigned in my other
courses, like INFO: 601 Foundations of Information and INFO: 657
Digital Humanities, readings on Queer Studies generate much interest
and discussion, with students often claiming them as their favorite
readings from the course. This course would pair well with another
Special Topics currently being proposed, "Disability Studies for
LIS." To better cater to the LIS student population, whom I believe
will be the ones most interested in the course, I would plan to have
incorporate direct research in a Queer archive in New York City.
\end{itemize}

\subsection{bank}
\label{sec:orga0ef5ce}
\begin{itemize}
\item active \& experiential learning: students practice hands-on with
knowledge they bring from experience
\begin{itemize}
\item student investment in material: students must have a stake in what
they are learning.
\item meeting the student where they are at: not making assumptions
about prior knowledge.
\end{itemize}
\end{itemize}
\end{document}
